\centerline{\bf \Huge ДЗ$\cdot$Делимость}
\title{ДЗ$\cdot$Делимость}

\begin{flushright}
    \scriptsize Ты поймёшь, как важно идти вперёд, \\как только сделаешь первый шаг.\\
    Мисато Кацураги. Евангелион.
\end{flushright}

\problem{1}В честь открытия ЛМШ работники решили устроить конфетный пир. Работники решили разыграть директора и договорились принести конфеты следующим образом. Первый работник положил в огромную пиалу одну конфету. А каждый последующий работник клал на одну конфету больше чем предыдущий. Всего работников было 2023, сможет ли директор ЛМШ поделить конфеты поровну между лососями и шиншиллами?

\problem{2}Баир решил доказать всем свои арифметические навыки и вручную начал вычислять $999^{999}$ и вычислил его абсолютно верно! Но Анджур Аркадьевич дабы удостовериться в умениях Баира попросил его вычислить сумму цифр этого числа. А потом у полученного числа еще раз вычислить сумму цифр, а потом еще раз, еще раз и так до тех пор пока не осталась цифра 9. Правильно ли вычислил цифру Баир?

\problem{3}Алина, Эрдем, Баир и Расул Владимирович приехали в Сириус и решили устроить турнир по поеданию пряников по следующим правилам: \\

\noindent1) Пряники поедаются по кругу.\\
2) Ход $-$ поедание пряников одним человеком.\\
3) Суммарное количество съеденных за последние 4 хода пряников делится поровну между 10 лососями, а то что осталось надо съесть в следующий ход.\\

Известно что в первый ход Алина съела 1 пряник, во второй ход Эрдем съел 2 пряника, в третий ход Баир съел 3 пряника, а в четвёртый ход Расул Владимирович съел аж 4 пряника(то есть дальше Алина съест 0 прянников, Эрдем 9, Баир 6 и так далее). Может ли случиться так, что в какой-то момент игры в три хода подряд было съедено по 9 пряников?

За количество пряников можете не волноваться, директор заранее в рюкзаке вынес из столовой бесконечное число пряников!